% Part: first-order-logic
% Chapter: completeness
% Section: henkin-expansion

\documentclass[../../../include/open-logic-section]{subfiles}

\begin{document}

\olfileid{fol}{com}{hen}

\olsection{De Henkin-uitbreiding}

\begin{explain}
Een lastig deel van het volledigheidsbewijs is dit: het model dat we uit
een volledige consistente verzameling~$\Gamma$ bouwen, moet alle
gekwantificeerde !!{formula}s in~$\Gamma$ waar maken. Een truc van Leon
Henkin zorgt daarvoor. We breiden eerst de taal uit met oneindig veel
nieuwe !!{constant}s. Daarna voegen we voor elke !!{formula} met één
vrije !!{variable} $!A(x)$ een formule toe van de vorm
\iftag{prvEx}
      {$\lexists[x][!A(x)] \lif !A(c)$}
      {$\lnot\lforall[x][!A(x)] \lif \lnot !A(c)$},
waarbij $c$ een van de nieuwe !!{constant}s is. Later bouwen we de
!!{structure} die~$\Gamma$ waar maakt. De toegevoegde formules zorgen er
dan voor dat elke
\iftag{prvEx}
{ware existentiële zin een getuige heeft}
{onware universele zin een tegenvoorbeeld heeft}
onder de nieuwe constanten.
\end{explain}

\begin{prop}
\ollabel{prop:lang-exp}
Stel dat $\Gamma$ consistent is in $\Lang L$. Stel ook dat we $\Lang L'$
uit $\Lang L$ maken door !!a{denumerable} verzameling nieuwe
!!{constant}s $\Obj d_0$, $\Obj d_1$, \dots toe te voegen. Dan blijft
$\Gamma$ consistent in~$\Lang L'$.
\end{prop}

\begin{defn}[Verzadigde verzameling]
Een verzameling $\Gamma$ van !!{formula}s van een taal $\Lang {L}$ heet
\emph{verzadigd} precies als aan de volgende voorwaarde is voldaan. Neem
een willekeurige !!{formula}~$!A(x) \in \Frm[L]$ met één vrije
!!{variable}~$x$. Er moet dan een !!{constant}~$c \in \Lang{L}$ zijn
waarvoor
\iftag{prvEx}
      {$\lexists[x][!A(x)] \lif !A(c) \in \Gamma$}
      {$\lnot\lforall[x][!A(x)] \lif \lnot !A(c) \in \Gamma$}.
\end{defn}

We hebben de volgende definitie nodig in het bewijs van de volgende
stelling.

\begin{defn}
\ollabel{defn:henkin-exp}
Neem $\Lang L'$ uit \olref{prop:lang-exp}. Zet $!A_0(x_0)$,
$!A_1(x_1)$, \dots in één lijst. Dit is een opsomming van alle
!!{formula}s~$!A_i(x_i)$ van~$\Lang L'$ waarin één variabele ($x_i$)
vrij voorkomt. We definiëren nu de !!{sentence}s~$!D_n$ stap voor stap,
door inductie naar~$n$.

Laat $c_0$ de eerste !!{constant} onder de $\Obj d_i$ zijn die we aan
~$\Lang{L}$ hebben toegevoegd en die niet in~$!A_0(x_0)$ voorkomt. Neem
nu aan dat $!D_0$, \dots,~$!D_{n-1}$ al zijn gedefinieerd. Laat $c_n$
dan de eerste van de nieuwe !!{constant}s~$\Obj d_i$ zijn die niet in
$!D_0$, \dots,~$!D_{n-1}$ en ook niet in~$!A_n(x_n)$ voorkomt.

Definieer $!D_{n}$ als de volgende !!{formula}: \iftag{prvEx}
{$\lexists[x_{n}][!A_{n}(x_{n})] \lif
  !A_{n}(c_{n})$}{$\lnot\lforall[x_{n}][!A_{n}(x_{n})] \lif \lnot
  !A_{n}(c_{n})$}.
\end{defn}

\begin{lem}
\ollabel{lem:henkin}
Elke consistente verzameling~$\Gamma$ past in een grotere verzameling
$\Gamma'$ die zowel verzadigd als consistent is.
\end{lem}

\begin{proof}
Begin met een consistente verzameling zinnen~$\Gamma$ in een
taal~$\Lang{L}$. Voeg !!a{denumerable} verzameling nieuwe !!{constant}s
aan de taal toe. Zo krijgen we~$\Lang{L'}$. Volgens
\olref{prop:lang-exp} blijft $\Gamma$ consistent in deze ruimere taal.
Neem verder de formules $!D_i$ uit \olref{defn:henkin-exp}. Definieer
\begin{align*}
\Gamma_0 & = \Gamma \\
\Gamma_{n+1} & = \Gamma_n \cup \{!D_n \}
\end{align*}
Dus $\Gamma_{n+1} = \Gamma \cup \{ !D_0, \dots, !D_n \}$. Definieer
ook $\Gamma' = \bigcup_{n} \Gamma_n$. Door de manier waarop we de
formules hebben toegevoegd, is $\Gamma'$ verzadigd.

Stel dat $\Gamma'$ inconsistent is. Dan moet er een $n$ zijn waarvoor
$\Gamma_n$ inconsistent is (Opgave: leg uit waarom). Om te bewijzen dat
$\Gamma'$ consistent is, hoeven we dus alleen met inductie naar~$n$ aan
te tonen dat elke verzameling~$\Gamma_n$ consistent is.

De basisstap zegt dat $\Gamma_0 = \Gamma$ consistent is. Dat is precies
de aanname van de stelling. Voor de inductiestap nemen we aan dat
$\Gamma_{n}$ consistent is, maar dat $\Gamma_{n+1} =
\Gamma_n \cup \{!D_n\}$ inconsistent is. De formule $!D_n$ is
\iftag{prvEx}
{$\lexists[x_{n}][!A_{n}(x_n)] \lif !A_{n}(c_{n})$}
{$\lnot\lforall[x_{n}][!A_{n}(x_n)] \lif \lnot !A_{n}(c_{n})$}.
Hier is $!A_n(x_n)$ !!a{formula} van $\Lang{L'}$ waarin alleen de
variabele~$x_n$ vrij voorkomt. Kijk nu terug naar de keuze van~$c_n$ in
\olref{defn:henkin-exp}. Daaruit volgt dat $c_n$ niet in~$!A_n(x_n)$ en
ook niet in~$\Gamma_n$ voorkomt.

Als $\Gamma_n \cup \{!D_n\}$ inconsistent is, dan $\Gamma_n
\Proves \lnot !D_n$. Daaruit volgen de twee beweringen hieronder:
\iftag{prvEx}{
\[
\Gamma_n \Proves \lexists[x_n][!A_n(x_n)]
\qquad
\Gamma_n \Proves \lnot !A_n(c_n)
\]}{
\[
\Gamma_n \Proves \lnot\lforall[x_n][!A_n(x_n)]
\qquad
\Gamma_n \Proves !A_n(c_n)
\]}
De constante $c_n$ komt niet voor in
$\Gamma_n$ of in~$!A_n(x_n)$. Daarom mogen we
\tagrefs{prfAX/{fol:axd:qpr:thm:strong-generalization},prfSC/{fol:seq:qpr:thm:strong-generalization},prfND/{fol:ntd:qpr:thm:strong-generalization},prfTab/{fol:tab:qpr:thm:strong-generalization}}
toepassen. Uit \iftag{prvEx}
{$\Gamma_n \Proves \lnot !A_n(c_n)$}
{$\Gamma_n \Proves !A_n(c_n)$}
volgt dan
\iftag{prvEx}
{$\Gamma_n \Proves \lforall[x_n][\lnot !A_n(x_n)]$}
{$\Gamma_n \Proves \lforall[x_n][!A_n(x_n)]$}.
We hebben dus zowel
\iftag{prvEx}
{$\Gamma_n \Proves \lexists[x_n][!A_n(x_n)]$ als
$\Gamma_n \Proves \lforall[x_n][\lnot !A_n(x_n)]$}
{$\Gamma_n \Proves \lnot\lforall[x_n][!A_n(x_n)]$ als
$\Gamma_n \Proves \lforall[x_n][!A_n(x_n)]$}.
Daarmee is $\Gamma_n$ zelf inconsistent.
\iftag{prvEx}
{(Merk op dat
\iftag{prvAll}
{$\lforall[x_n][\lnot !A_n(x_n)] \Proves
  \lnot\lexists[x_n][!A_n(x_n)]$}
{$\lforall[x_n][\lnot !A_n]$ per definitie gelijk is aan
  $\lnot\lexists[x_n][\lnot\lnot !A_n(x_n)]$ en dat
  $\lnot\lexists[x_n][\lnot\lnot !A_n(x_n)] \Proves
  \lnot\lexists[x_n][!A_n(x_n)]$}.)}{}
Dat is een tegenspraak: we hadden aangenomen dat $\Gamma_n$ consistent
was. Dus is $\Gamma_n \cup \{ !D_n\}$ consistent.
\end{proof}

\begin{explain}
Nu bekijken we wat verzadiging ons oplevert. Neem een verzameling die
\emph{volledig}, consistent en verzadigd is. Dan geldt:
\iftag{prvAll}{zij bevat een universeel gekwantificeerde !!{sentence}
  precies als zij al haar instanties bevat}{}\iftag{defAll,defEx}{}{
  en }\iftag{prvAll}{zij bevat een existentieel gekwantificeerde
  !!{sentence} precies als zij minstens één instantie bevat}{}. We
gebruiken dit straks om te bewijzen dat de !!{structure} die we uit een
volledige, consistente, verzadigde verzameling bouwen al haar
gekwantificeerde zinnen waar maakt.
\end{explain}

\begin{prop}\ollabel{prop:saturated-instances}
Stel dat $\Gamma$ volledig, consistent en verzadigd is.
\begin{tagenumerate}{prvEx,prvAll}
\tagitem{prvEx}{$\lexists[x][!A(x)] \in \Gamma$ precies als
  $!A(t) \in \Gamma$ voor minstens één gesloten term~$t$.}{}
\tagitem{prvAll}{$\lforall[x][!A(x)] \in \Gamma$ precies als
  $!A(t) \in \Gamma$ voor alle gesloten termen~$t$.}{}
\end{tagenumerate}
\end{prop}

\begin{proof}
  \begin{tagenumerate}{prvEx,prvAll}
  \tagitem{prvEx}{%
    \iftag{probEx}{Opgave.}{Stel eerst dat $\lexists[x][!A(x)]
      \in \Gamma$. Omdat $\Gamma$ verzadigd is, geldt
      $(\lexists[x][!A(x)] \lif !A(c)) \in \Gamma$ voor een
      !!{constant}~$c$. Uit
      \tagrefs{prfAX/{fol:axd:ppr:prop:provability-lif},prfSC/{fol:seq:ppr:prop:provability-lif},prfND/{fol:ntd:ppr:prop:provability-lif},prfTab/{fol:tab:ppr:prop:provability-lif}}, onderdeel~(1),
      en \olref[ccs]{prop:ccs}\olref[ccs]{prop:ccs-prov-in} volgt dat
      $!A(c) \in \Gamma$.

    Voor de andere richting hebben we verzadiging niet nodig. Stel dat
    $!A(t) \in \Gamma$. Dan volgt uit
      \tagrefs{prfAX/{fol:axd:qpr:prop:provability-quantifiers},prfSC/{fol:seq:qpr:prop:provability-quantifiers},prfND/{fol:ntd:qpr:prop:provability-quantifiers},prfTab/{fol:tab:qpr:prop:provability-quantifiers}}, onderdeel~(1),
    dat $\Gamma \Proves \lexists[x][!A(x)]$. Volgens
    \olref[ccs]{prop:ccs}\olref[ccs]{prop:ccs-prov-in} geldt dan
    $\lexists[x][!A(x)] \in \Gamma$.}}{}
  \tagitem{prvAll}{%
    \iftag{probAll}{Opgave.}{Stel dat $!A(t) \in \Gamma$ voor alle
      gesloten termen~$t$. Neem voor een tegenspraak aan dat
      $\lforall[x][!A(x)] \notin \Gamma$. Omdat $\Gamma$ volledig is,
      geldt $\lnot\lforall[x][!A(x)] \in \Gamma$. Uit verzadiging volgt
      dat \iftag{prvEx}{$(\lexists[x][\lnot !A(x)] \lif \lnot !A(c)) \in
        \Gamma$}{$(\lnot\lforall[x][!A(x)] \lif \lnot !A(c)) \in
        \Gamma$} voor een !!{constant}~$c$. Volgens de aanname geldt,
      omdat $c$ een gesloten term is, $!A(c) \in \Gamma$. Maar dan zou
      $\Gamma$ inconsistent zijn. (Opgave: geef de !!{derivation} die
      laat zien dat
      \[
      \lnot \lforall[x][!A(x)],
      \iftag{prvEx}{\lexists[x][\lnot !A(x)] \lif \lnot
        !A(c)}{\lnot\lforall[x][!A(x)] \lif \lnot
        !A(c)}, !A(c)
      \]
      inconsistent is.)

      Voor de andere richting hebben we verzadiging niet nodig. Stel dat
      $\lforall[x][!A(x)] \in \Gamma$. Uit
      \tagrefs{prfAX/{fol:axd:qpr:prop:provability-quantifiers},prfSC/{fol:seq:qpr:prop:provability-quantifiers},prfND/{fol:ntd:qpr:prop:provability-quantifiers},prfTab/{fol:tab:qpr:prop:provability-quantifiers}},
      onderdeel~(2), volgt dan $\Gamma \Proves !A(t)$. Volgens
      \olref[ccs]{prop:ccs} geldt $!A(t) \in \Gamma$.}}{}
  \end{tagenumerate}
\end{proof}

\end{document}
